\documentclass[11pt]{article}
\usepackage[a4paper,margin=1in]{geometry}
\usepackage[T1]{fontenc}
\usepackage{microtype}
\usepackage{amsmath,amssymb,amsthm,mathtools}
\usepackage{enumitem}
\usepackage{booktabs}
\usepackage{xcolor}
\usepackage{tikz}
\PassOptionsToPackage{hyphens}{url}
\usepackage[hidelinks]{hyperref}

\setlength{\parskip}{0.45em}
\setlength{\parindent}{0pt}

\theoremstyle{plain}
\newtheorem{theorem}{Theorem}[section]
\newtheorem{lemma}[theorem]{Lemma}
\newtheorem{proposition}[theorem]{Proposition}
\newtheorem{corollary}[theorem]{Corollary}
\newtheorem{citedtheorem}[theorem]{Cited Theorem}
\newtheorem{maintheorem}{Theorem}
\renewcommand{\themaintheorem}{\Alph{maintheorem}}
\theoremstyle{definition}
\newtheorem{definition}[theorem]{Definition}
\newtheorem{computation}[theorem]{Computation}
\newtheorem{problem}[theorem]{Open Problem}
\theoremstyle{remark}
\newtheorem{remark}[theorem]{Remark}

\newcommand{\T}{\mathbin{\square}}
\newcommand{\Z}{\mathbb{Z}}
\newcommand{\why}[1]{\unskip\ \textit{[#1]}}
\newenvironment{steps}{\begin{enumerate}[label=(\arabic*),leftmargin=2.2em,itemsep=0.35em]}{\end{enumerate}}

\title{The Toroidal Grid $C_3 \T C_n$ Is 3-Choosable\\for Every $n$}
\author{Amrutanshu \and Siddharth Narayan Mishra}
\date{Draft, 28 October 2026}

\begin{document}
\maketitle

\begin{abstract}
Cai, Wang and Zhu conjectured that every toroidal grid $C_m \T C_n$ is 3-choosable. Li, Shao,
Petrov and Gordeev proved this when $m$ or $n$ is even, and showed that the Alon--Tarsi number
of $C_m \T C_n$ is $4$ when $m$ and $n$ are both odd, so that the Alon--Tarsi method cannot
reach the remaining case. We prove that $C_3 \T C_n$ is 3-choosable for every odd $n$. Hence
$\chi_\ell(C_3 \T C_n) = 3$ for every $n \ge 3$. The proof cuts the torus open along two columns,
colours the remaining strip with the Combinatorial Nullstellensatz while pinning one colour and
forbidding one other, and then fills the two columns by hand; the forbidden colour is chosen in
advance so that the fill cannot fail. We also record a second, kernel-method proof for the
family $n = 3k$, $k$ odd, and what is known so far for $C_m \T C_n$ with $m, n \ge 5$ odd.
\end{abstract}

\begin{center}
\fbox{\begin{minipage}{0.93\textwidth}\small
\textbf{Status of this draft.} Every step of the proof of Theorem~A is either proved in full
here or is a finite statement checked by exhaustive computation over all list assignments up to
renaming of colours (Lemma~\ref{lem:prism}). The scripts are listed in
Section~\ref{sec:verification}. Still to do before submission: (i) a hand proof of
Lemma~\ref{lem:prism} (Amrutanshu's prism lemma has one, by Hall's theorem, to be merged);
(ii) a full literature check of the papers citing [LSPG] and of the dynamic survey [S], to
confirm that the odd case of $C_3 \T C_n$ is not already settled.
\end{minipage}}
\end{center}

%=====================================================================================
\section{Introduction}

All graphs are finite and simple. A \emph{list assignment} $L$ of a graph $G$ gives each vertex
$v$ a set $L(v)$ of colours. An \emph{$L$-colouring} is a proper colouring $c$ with $c(v) \in L(v)$
for every $v$. The graph $G$ is \emph{$k$-choosable} if it has an $L$-colouring whenever
$|L(v)| \ge k$ for all $v$, and the \emph{choice number} $\chi_\ell(G)$ is the least such $k$.

For $m, n \ge 3$ the \emph{toroidal grid} $C_m \T C_n$ has vertex set $\Z_m \times \Z_n$, and
$(i, j)$ is adjacent to $(i \pm 1, j)$ and to $(i, j \pm 1)$. We call $i$ the \emph{row} and $j$
the \emph{column}. Column $j$ induces a cycle of length $m$; row $i$ induces a cycle of length
$n$. When $m = 3$ every column is a triangle, so $\chi_\ell(C_3 \T C_n) \ge 3$.

\begin{citedtheorem}[{Li, Shao, Petrov and Gordeev [LSPG]}]\label{thm:lspg}
If $m$ or $n$ is even then $C_m \T C_n$ is 3-choosable. If $m$ and $n$ are both odd then the
Alon--Tarsi number of $C_m \T C_n$ equals $4$.
\end{citedtheorem}

So for odd $m, n$ the conjecture of Cai, Wang and Zhu [CWZ] is open, and the polynomial method
on its own cannot settle it. The smallest case $C_3 \T C_3 = K_3 \T K_3$ is the line graph of
$K_{3,3}$, and Galvin's theorem [G] shows that it is 3-choosable.

\begin{maintheorem}\label{thm:main}
For every odd $n \ge 3$ the toroidal grid $C_3 \T C_n$ is 3-choosable. Consequently
$\chi_\ell(C_3 \T C_n) = 3$ for every $n \ge 3$.
\end{maintheorem}

The second sentence follows from the first, from Cited Theorem~\ref{thm:lspg} for even $n$, and
from the triangles. The case $n = 3$ is Galvin's theorem. For odd $n \ge 5$ the proof is in
Section~\ref{sec:proof}. It uses two tools: an \emph{odd strip lemma} (Section~\ref{sec:strip}),
proved with the Combinatorial Nullstellensatz, and a \emph{two-column lemma} for the triangular
prism (Section~\ref{sec:prism}).

%=====================================================================================
\section{Tools from the polynomial method}

For a graph $G$ with an orientation of its edges, the \emph{graph polynomial} is
$P_G = \prod_{(u,v)} (x_u - x_v)$, the product over all arcs.

\begin{citedtheorem}[Combinatorial Nullstellensatz, Alon {[A]}]\label{thm:cn}
Let $P \in \mathbb{F}[x_1, \dots, x_N]$ have degree $\sum_i t_i$, and suppose the coefficient of
$\prod_i x_i^{t_i}$ in $P$ is nonzero. If $S_1, \dots, S_N \subseteq \mathbb{F}$ with
$|S_i| > t_i$, then some $s \in S_1 \times \dots \times S_N$ has $P(s) \ne 0$.
\end{citedtheorem}

Applied to $P_G$ with $|L(v)| \ge t_v + 1$ (colours viewed as complex numbers), a point where
$P_G \ne 0$ is an $L$-colouring.

\begin{citedtheorem}[Coefficient formula, Lasoń {[La]}, Karasev and Petrov {[KP]}]\label{thm:coef}
Let $P$ have degree $\sum_i t_i$ and let $A_i \subseteq \mathbb{F}$ with $|A_i| = t_i + 1$. Put
$\varphi_i(z) = \prod_{b \in A_i} (z - b)$. Then
\[
[x_1^{t_1} \cdots x_N^{t_N}]\, P \;=\; \sum_{a \in A_1 \times \dots \times A_N}
\frac{P(a)}{\varphi_1'(a_1) \cdots \varphi_N'(a_N)} .
\]
\end{citedtheorem}

Throughout, $\omega = e^{2\pi i/3}$ and
\[
g \;=\; 3(\omega - \omega^2) \;=\; 3 i \sqrt3, \qquad g^2 = -27 .
\]
A direct expansion gives $(1-\omega)^3 = -3\omega + 3\omega^2 = -g$ and
$(1-\omega^2)^3 = 3\omega - 3\omega^2 = g$. Define $\chi : \Z_3 \to \{0, \pm 1\}$ by
$\chi(0) = 0$, $\chi(1) = -1$, $\chi(2) = 1$. Then
\begin{equation}\label{eq:chi}
(1 - \omega^d)^3 = g\,\chi(d) \qquad (d \in \Z_3).
\end{equation}

%=====================================================================================
\section{The odd strip lemma}\label{sec:strip}

For $\ell \ge 2$ let $H_\ell = C_3 \T P_\ell$, with vertex set $\Z_3 \times \{1, \dots, \ell\}$:
each column is a triangle and row $i$ is the path $(i,1), \dots, (i,\ell)$. It has $3\ell$
vertices, $3\ell$ column edges and $3(\ell - 1)$ row edges, so $6\ell - 3$ edges.

An \emph{end weight} is a map $\mu : V(H_\ell) \to \{0, 1, 2\}$, zero outside columns $1$ and
$\ell$, with $\sum_v \mu(v) = 3$. Put
\[
k = \sum_{i} \mu(i, \ell), \qquad s = \sum_{v} \mu(v)\, \mathrm{row}(v) \bmod 3 .
\]

\begin{lemma}[Odd strip lemma]\label{lem:strip}
Let $\ell \ge 3$ be odd and let $\mu$ be an end weight with $k \in \{1, 2\}$ and $s \not\equiv 0
\pmod 3$. Then $H_\ell$ has an $L$-colouring for every list assignment with
$|L(v)| \ge 3 - \mu(v)$ for all $v$.
\end{lemma}

So one may \emph{pin} a vertex of weight $2$ (list of size $1$) and \emph{forbid} one colour at a
vertex of weight $1$ (list of size $2$).

\begin{proof}
Orient each column edge $(i, c) \to (i+1, c)$ and each row edge $(i, c) \to (i, c+1)$, and let
$P$ be the graph polynomial. Put $t_v = 2 - \mu(v) \ge 0$. Then $\sum_v t_v = 6\ell - 3 = \deg P$,
so by Cited Theorem~\ref{thm:cn} it suffices to show that
$c_\mu := [\prod_v x_v^{t_v}]\, P \ne 0$.

\begin{steps}
\item \emph{Reduction to cube roots.} Let $Q = P \cdot \prod_v x_v^{\mu(v)}$. Then
$\deg Q = 6\ell$ and $c_\mu = [\prod_v x_v^2]\, Q$. Apply Cited Theorem~\ref{thm:coef} to $Q$
with every $t_v = 2$ and every $A_v = \{1, \omega, \omega^2\}$. Here $\varphi(z) = z^3 - 1$ and
$\varphi'(a) = 3a^2 = 3a^{-1}$ for $a^3 = 1$. Hence
\[
c_\mu = 3^{-3\ell} \sum_{\phi} P(\phi) \prod_v \phi(v)^{1 + \mu(v)},
\]
summed over all maps $\phi : V(H_\ell) \to \{1, \omega, \omega^2\}$.

\item \emph{Column states.} $P(\phi) = 0$ unless $\phi$ is proper, and in particular unless
$\phi$ is injective on every column. The bijections $\Z_3 \to \Z_3$ are exactly the six maps
$i \mapsto a + \varepsilon i$ with $a \in \Z_3$, $\varepsilon = \pm 1$. So the contributing
$\phi$ have the form $\phi(i, c) = \omega^{a_c + \varepsilon_c i}$.

\item \emph{Column edges.} For one column,
$\prod_{i \in \Z_3} (\omega^{a + \varepsilon i} - \omega^{a + \varepsilon (i+1)})
= \omega^{3a + 3\varepsilon}(1 - \omega^\varepsilon)^3 = (1 - \omega^{\varepsilon})^3
= -\varepsilon g$.

\item \emph{Row edges.} Take columns $c, c+1$ with states $(a, \varepsilon)$ and
$(b, \varepsilon')$. If $\varepsilon' = -\varepsilon$, the colours of $(i, c)$ and $(i, c+1)$
agree when $a + \varepsilon i \equiv b - \varepsilon i$, that is, for $i \equiv 2\varepsilon (b -
a)$. So some row edge has equal ends and the product is $0$. If $\varepsilon' = \varepsilon$,
the product over the three row edges is
$\prod_i \omega^{a + \varepsilon i}(1 - \omega^{b - a}) = (1 - \omega^{b-a})^3
= g\,\chi(b - a)$ by \eqref{eq:chi}. Hence $\varepsilon_c = \varepsilon$ is constant, and with
$d_c = a_{c+1} - a_c$,
\[
P(\phi) = (-\varepsilon g)^{\ell} \, g^{\ell - 1} \prod_{c=1}^{\ell - 1} \chi(d_c).
\]

\item \emph{Vertex weights.} In column $c$, $\sum_i (a_c + \varepsilon i) = 3a_c + 3\varepsilon
\equiv 0$, and $\sum_i (a_c + \varepsilon i)\mu(i,c) = a_c M_c + \varepsilon R_c$ with
$M_c = \sum_i \mu(i,c)$ and $R_c = \sum_i i\,\mu(i,c)$. Now $M_1 = 3 - k$, $M_\ell = k$, the other
$M_c$ vanish, and $R_1 + R_\ell \equiv s$. Since $3a_1 \equiv 0$,
\[
\prod_v \phi(v)^{1 + \mu(v)} = \omega^{k (a_\ell - a_1) + \varepsilon s}.
\]

\item \emph{Summation.} The map $(a_1, \dots, a_\ell) \mapsto (a_1, d_1, \dots, d_{\ell-1})$ is a
bijection and $a_\ell - a_1 = \sum_c d_c$. So
\[
c_\mu = 3^{-3\ell}\, g^{2\ell - 1} \Big(\sum_{\varepsilon = \pm 1} (-\varepsilon)^{\ell}
\omega^{\varepsilon s}\Big) \cdot 3 \Big(\sum_{d \in \Z_3} \chi(d)\,\omega^{kd}\Big)^{\ell - 1}.
\]
For odd $\ell$ the first bracket is $\omega^{-s} - \omega^{s} = \omega^{2s} - \omega^{s}$, which is
nonzero because $s \not\equiv 0$. The second bracket is $\mu_k := \omega^{2k} - \omega^{k}$,
which is nonzero because $k \not\equiv 0 \pmod 3$. Since $g \ne 0$, $c_\mu \ne 0$. \qedhere
\end{steps}
\end{proof}

\begin{remark}
Simplifying with $g^2 = -27$ gives $c_\mu = \tfrac{g}{9}\, \mu_k^{\ell - 1}(\omega^{2s} -
\omega^{s}) = \pm 3^{(\ell-1)/2}$ for odd $\ell$, $k \in \{1,2\}$, $s \not\equiv 0$. This closed
form (for every $\ell \ge 2$ and every end weight) was checked against three independent exact
computations for $\ell$ from $2$ to $11$, with no mismatch (script \texttt{ws-W1}). The lemma and
the computation first appear in Amrutanshu's notes; the proof above is the one we re-derived.
\end{remark}

%=====================================================================================
\section{The two-column lemma}\label{sec:prism}

The \emph{prism} $\Pi = K_3 \T K_2$ has two triangles $u_0u_1u_2$ and $v_0v_1v_2$ and the three
edges $u_tv_t$.

\begin{lemma}[Two-column lemma]\label{lem:prism}
Let $A_0, A_1, A_2, B_0, B_1, B_2$ be lists for $u_0, u_1, u_2, v_0, v_1, v_2$ with $|A_0| = 3$
and every other list of size $2$ or $3$. If $\Pi$ has no colouring from these lists, then
$|B_1| = |B_2| = 2$ and $B_1 = B_2$.
\end{lemma}

\begin{proof}[Proof (exhaustive computation)]
Colourability does not change when colours are renamed. Read the lists in the order $A_0, A_1,
A_2, B_0, B_1, B_2$ and rename colours in order of first appearance. Afterwards $A_0 = \{0, 1,
2\}$, and each later list consists of some colours already used together with an initial segment
$u, u+1, \dots$ of the unused ones, where $u$ is the number of colours used so far. The program
\texttt{prism\_sizes.c} enumerates all $236\,160\,170$ assignments of this normal form, decides
colourability of each by complete search, and finds $8\,335$ that are not colourable. All of them
have $|B_1| = |B_2| = 2$ and $B_1 = B_2$.
\end{proof}

\begin{remark}[Amrutanshu's prism lemma]
For sizes exactly $(3,2,2,2,2,2)$ the non-colourable assignments are exactly those of the two
forms (Z) $B_0 = B_1 = B_2$, and (W) $A_1 = A_2 = D$, $A_0 = D \cup \{\zeta\}$, $B_1 = B_2 = E$,
$B_0 \setminus E = \{\zeta\}$. In (W), $u_1, u_2$ use up $D$ and force $u_0 = \zeta$, while
$v_1, v_2$ use up $E$ and force $v_0 = \zeta$. Amrutanshu proves this by hand with Hall's
theorem; \texttt{prism.c} confirms it on all $397\,546$ normal-form assignments ($986$ not
colourable, no exception). Both (Z) and (W) have $B_1 = B_2$, which is all the proof uses.
\end{remark}

%=====================================================================================
\section{Proof of Theorem A}\label{sec:proof}

Let $n \ge 5$ be odd and let $L$ be a list assignment of $T = C_3 \T C_n$ with $|L(v)| = 3$ for
every $v$ (larger lists can be shrunk). We consider two cases.

\subsection*{Case 1: some row list is not contained in the next one}

Suppose there are $r \in \Z_3$, $j \in \Z_n$ and a colour
\[
\alpha \in L(r, j) \setminus L(r, j+1).
\]
Write $C_i$ for column $j + i$, and let $r', r''$ be the two rows other than $r$.

\begin{steps}
\item \emph{Choice of $\beta$.} If $|L(r', C_2) \cap L(r'', C_2)| = 2$, let $\beta$ be the unique
colour of $L(r', C_2) \setminus L(r'', C_2)$. Otherwise let $\beta$ be any colour.

\item \emph{The strip.} Delete columns $C_1$ and $C_2$. Since $n \ge 5$, what remains is a copy
of $H_{n-2}$ with columns $C_3, C_4, \dots, C_{n-1}, C_0$ in this order, and $n - 2 \ge 3$ is odd.
Let $\mu(r, C_0) = 2$, $\mu(r', C_3) = 1$, and $\mu = 0$ elsewhere. The last column is $C_0$, so
$k = 2$, and $s \equiv 2r + r' \equiv r' - r \not\equiv 0$.

\item \emph{Colour the strip.} Give $(r, C_0)$ the list $\{\alpha\}$, give $(r', C_3)$ the list
$L(r', C_3) \setminus \{\beta\}$, and keep every other list. These satisfy $|L(v)| \ge 3 -
\mu(v)$, so Lemma~\ref{lem:strip} gives an $L$-colouring $\psi$ of the strip with
$\psi(r, C_0) = \alpha$ and $\psi(r', C_3) \ne \beta$.

\item \emph{The lists left for $C_1 \cup C_2$.} The subgraph on $C_1 \cup C_2$ is a prism with
$u_t = (t, C_1)$ and $v_t = (t, C_2)$. A vertex of $C_1$ has exactly one neighbour outside
$C_1 \cup C_2$, namely its row neighbour in $C_0$; a vertex of $C_2$ has exactly one, its row
neighbour in $C_3$. (These are different columns since $n \ge 5$.) So $\psi$ extends to $T$ if and
only if the prism is colourable from
\[
A_t = L(t, C_1) \setminus \{\psi(t, C_0)\}, \qquad B_t = L(t, C_2) \setminus \{\psi(t, C_3)\}.
\]
All these lists have size $2$ or $3$, and $A_r = L(r, C_1)$ has size $3$ because
$\psi(r, C_0) = \alpha \notin L(r, C_1)$.

\item \emph{The fill cannot fail.} Suppose it does. By Lemma~\ref{lem:prism} (with $u_0 = (r,
C_1)$ and $\{u_1, u_2\} = \{(r', C_1), (r'', C_1)\}$), $B_{r'} = B_{r''} = E$ for a 2-set $E$.
Then $y' := \psi(r', C_3) \in L(r', C_2)$ and $y'' := \psi(r'', C_3) \in L(r'', C_2)$, so
$L(r', C_2) = E \cup \{y'\}$ and $L(r'', C_2) = E \cup \{y''\}$, with $y', y'' \notin E$. Also
$y' \ne y''$, since $(r', C_3)$ and $(r'', C_3)$ are adjacent. Hence $L(r', C_2) \cap L(r'',
C_2) = E$ has size $2$, and $y'$ is the unique colour of $L(r', C_2) \setminus L(r'', C_2)$. By
step (1) this colour is $\beta$, so $\psi(r', C_3) = \beta$, contradicting step (3).
\end{steps}

So $\psi$ extends to an $L$-colouring of $T$.

\subsection*{Case 2: every row list is contained in the next one}

Then $L(r, j) \subseteq L(r, j+1)$ for all $r$ and $j$. All lists have size $3$, so $L(r, j) =
L(r, j+1)$, and each row has a constant list $A_r$. By Galvin's theorem, $C_3 \T C_3 =
L(K_{3,3})$ is 3-choosable, so it has a colouring in which every vertex of row $r$ gets a colour
from $A_r$. Let $X, Y, Z : \Z_3 \to$ colours be its three columns. Each is a proper colouring of a
triangle with $X(r), Y(r), Z(r) \in A_r$, and $X(r), Y(r), Z(r)$ are pairwise distinct, since row
$r$ of $C_3 \T C_3$ is a triangle. Colour column $j$ of $T$ by $X$ if $j = 0$, by $Y$ if $j$ is
odd and by $Z$ if $j \ge 2$ is even. Consecutive columns get $(X, Y)$, $(Y, Z)$ or $(Z, Y)$, and
the pair $(n-1, 0)$ gets $(Z, X)$ because $n - 1$ is even. All these pairs are compatible, so the
colouring is proper.

The two cases cover every list assignment, which proves Theorem~A. \qed

\begin{remark}[Credit and history]
The route (edge deletion, the strip lemma, the prism lemma, and a two-case proof) is due to
Amrutanshu. His first version split cases on $F(j) = \bigcap_r L(r, j)$ and needed a
propagation lemma, $L(r, j) \setminus F(j) \subseteq L(r, j+1)$ for bad assignments, to reach
Case~2. Case~1 above only needs $\alpha \in L(r, j) \setminus L(r, j+1)$, so splitting on that
directly makes the propagation lemma (and with it the edge-deletion theorem) unnecessary. This
simplification is in Amrutanshu's work plan of 26 October. We checked it by hand and, before
Lemma~\ref{lem:prism} was verified exhaustively, by testing the whole extension step on every
list assignment of $C_1 \cup C_2$ over palettes of $5$ and $6$ colours ($200\,000$ and
$6\,400\,000$ cases, no failure; with the choice of $\beta$ removed, $14\,592$ failures over $5$
colours).
\end{remark}

%=====================================================================================
\section{A second proof for $n = 3k$, $k$ odd: the kernel method}\label{sec:kernel}

For $n = 3k$ there is a proof of a different kind, by Amrutanshu, which we checked line by
line. Colour $C_3 \T C_{3k}$ by $\varphi(i,j) = (i + j) \bmod 3$ and orient column edges towards
the smaller colour and row edges towards the larger colour. Every vertex then has out-degree
exactly $2$. A minimal-counterexample argument shows that this orientation is
\emph{kernel-perfect} when $k$ is odd, and the kernel lemma of Bondy, Boppana and Siegel gives
3-choosability. Oddness of $k$ is used once, to walk around the columns in steps of $2$.

Our computations (workstream W5) agree: this orientation is kernel-perfect exactly when
$k \not\equiv 2 \pmod 4$. For $k \equiv 2 \pmod 4$ the orientation contains an induced directed
cycle of odd length $3n/2$. For $n = 5, 7, 11, 13$ no orientation of $C_3 \T C_n$ with all
out-degrees $2$ is kernel-perfect, so the kernel method is limited to $3 \mid n$. The kernel lemma
also holds for paintability, so this proof gives the stronger statement that $C_3 \T C_{3k}$ is
3-paintable for odd $k$. The proof of Theorem~A does not give paintability, because $\beta$ is
chosen from the lists in advance.

%=====================================================================================
\section{Towards $C_m \T C_n$ with $m, n \ge 5$ odd}\label{sec:general}

\subsection{What carries over}

\begin{corollary}\label{cor:rowconst}
Let $m, n \ge 3$ be odd.
\begin{enumerate}[label=(\alph*),nosep]
\item $C_m \T C_3$ is 3-choosable.
\item If $L$ is a 3-list assignment of $C_m \T C_n$ that is constant along every row, or constant
along every column, then $C_m \T C_n$ is $L$-colourable.
\end{enumerate}
\end{corollary}

\begin{proof}
(a) $C_m \T C_3 \cong C_3 \T C_m$, and Theorem~A applies. (b) Say $L(i, j) = A_i$. By (a),
$C_m \T C_3$ has a colouring with colours from $A_i$ in row $i$; its three columns $X, Y, Z$ are
proper colourings of the $m$-cycle, pairwise compatible. Colour the columns of $C_m \T C_n$ by
$X, Y, Z, Y, Z, \dots, Y, Z$ as in Case~2. Columns are symmetric.
\end{proof}

\begin{theorem}[Edge deletion for $m = 5, 7$; computer-assisted, W4]\label{thm:edge}
Let $m \in \{5, 7\}$ and $n \ge 3$ odd. For every edge $e$ of $G = C_m \T C_n$ and every vertex $w$
not on the cycle of $G$ (column or row) through $e$, the coefficient of
$x_w \prod_{v \ne w} x_v^2$ in $P_{G - e}$ is nonzero. Hence if some edge $uv$ has $L(u) \cap L(v)
= \emptyset$, then $G$ is $L$-colourable.
\end{theorem}

The proof uses the $30 \times 30$ (for $m = 5$) column transfer matrix: it is diagonalisable,
a unique pair $\pm i\lambda_1$ ($\lambda_1 \approx 5.6232$) dominates, the dominant amplitudes are
exactly nonzero, and the finitely many small $n$ are computed exactly. The case $m = 3$ is
Amrutanshu's. So in a bad 3-list assignment, every edge has overlapping lists.

\begin{computation}[W2]\label{comp:c5c5}
There is no list assignment of $C_5 \T C_5$ with lists of size $3$ from a palette of $4$ colours
that has no colouring. (Counterexample-guided search; the final master problem is certified
unsatisfiable by an independent solver, and every learned clause and each of its $4\,800$
symmetry images is re-verified. The same holds for $C_3 \T C_5$, as it must by Theorem~A.)
\end{computation}

\subsection{The obstacle: two-column extension for 5-cycles}

The proof of Theorem~A generalises step by step except one. Pick $\alpha \in L(r, j) \setminus
L(r, j+1)$, delete columns $j+1, j+2$, and colour the strip $C_m \T P_{n-2}$. The strip has $mk$
vertices and $2mk - m$ edges ($k = n - 2$), so the Nullstellensatz leaves $m$ units of freedom at
its ends: $2$ to pin $(r, j)$ to $\alpha$, and $1$ to forbid a colour $\beta_t$ at each row $t$ of
column $j + 3$ outside a pair $\{p, q\}$.

\emph{Strip coefficients.} Amrutanshu observed, for $m \in \{3,5,7,9\}$ and $k \le 9$, that the
coefficient of this shape vanishes exactly when $k$ is odd and $2r \equiv p + q \pmod m$. The
zeros are forced: the row reflection fixing $r$ swaps $p$ and $q$, fixes the shape, and reverses
the $mk$ column edges, an odd number. We confirmed the rule exactly for $m = 5$ and every strip
length from $2$ to $27$ ($2\,600$ cases, no exception).

\emph{The extension.} What remains is to fill $C_m \T K_2$ (the two deleted columns). For $m = 3$
this is Lemma~\ref{lem:prism}. For $m = 5$ we model a \emph{window}: lists of size $3$ on the two
deleted columns, the pin at row $0$, and an adversary who chooses any proper colouring of the
$5$-cycles in columns $j$ and $j+3$ (with $(0, j)$ coloured $\alpha$). A \emph{shape} is an end
weight with $c_S \ne 0$; it lets us forbid one colour at chosen rows. A window is \emph{good} if
some shape and some forbidden colours make the fill succeed against every adversary.

\begin{computation}[W8, $m = 5$, palette of $4$ colours]\label{comp:windows}
\begin{enumerate}[label=(\alph*),nosep]
\item The plan's shape (forbid at the rows of column $j+3$ outside $\{p, q\}$, $2r \not\equiv p+q$)
fails on $6\,904$ of the $4^{10} = 1\,048\,576$ windows. Random windows from $5$ or $6$ colours
($60\,000$) never fail.
\item There are $74$ end weights with the pin at row $0$, all other weights $0$ or $1$, and
$c_S \ne 0$ for \emph{every} odd strip length from $3$ to $27$; $66$ of them forbid colours on the
$\alpha$ side. Allowing all $74$ leaves $400$ bad windows, and $184$ once $\alpha$ is taken to be
the actual colour (with $4$ colours it is forced).
\item A cut can be made forwards (delete $j+1, j+2$) or backwards (delete $j, j-1$, using
$\alpha' \in L(r, j+1) \setminus L(r, j)$), at any row $r$ and column $j$ with $L(r, j) \ne L(r,
j+1)$. For $n \in \{5, 7\}$, every 3-list assignment of $C_5 \T C_n$ from $4$ colours in which
\emph{every} cut is bad is row-constant, column-constant, or has four rows with one common
constant list: there are exactly $600$ (for $n = 5$) and $2\,520$ (for $n = 7$) assignments of
the last kind, and no others. (For $n = 9$ we have only checked that such assignments exist.)
\end{enumerate}
\end{computation}

The last kind is coloured directly.

\begin{lemma}[One free row]\label{lem:freerow}
Let $m \ge 5$ be odd and $n \ge 3$. If $L$ is a 3-list assignment of $C_m \T C_n$ and there are a
row $i_0$ and a set $A$ with $L(i, j) = A$ for all $i \ne i_0$ and all $j$, then $C_m \T C_n$ is
$L$-colourable.
\end{lemma}

\begin{proof}
The rows other than $i_0$ form a path $R = (i_0 + 1, \dots, i_0 - 1)$ of $m - 1 \ge 4$ rows, an even
number. Choose a proper colouring $p : R \to \Z_3$ of this path whose two ends get the same value:
$0, 1, 2, 0$ followed by $1, 0, 1, 0, \dots$ for longer paths. Let $\rho(x) = x + 1$ on $\Z_3$.
Choose $\sigma_0, \dots, \sigma_{n-1} \in \{\rho, \rho^{-1}\}$ with $\sigma_0 \cdots
\sigma_{n-1} = \mathrm{id}$; this needs $k$ copies of $\rho$ with $2k \equiv n \pmod 3$, and $k \in
\{0, 1, 2\}$ works. Fix a bijection $\pi_0 : \Z_3 \to A$, put $\pi_{j+1} = \pi_j \sigma_j$, and
colour $(i, j)$ with $\pi_j(p(i))$ for $i \in R$. Columns are proper because $p$ is, and row edges
are proper because $\sigma_j$ has no fixed point. The two neighbours of $(i_0, j)$ in column $j$
get the same colour $a_j := \pi_j(p(\text{end}))$, so $(i_0, j)$ keeps at least $2$ colours of
$L(i_0, j) \setminus \{a_j\}$, and row $i_0$ is a cycle with lists of size at least $2$.

Such a cycle is colourable unless it is odd and all its lists are one 2-set $B$ [ERT]. Suppose
this happens. Then $L(i_0, j) = B \cup \{a_j\}$ with $a_j \notin B$ for every $j$, and since
$a_{j+1} \ne a_j$, the colours $a_j$ take at least two values in $A \setminus B$. Now replace every
$\pi_j$ by $\tau \pi_j$ for a 3-cycle $\tau$ of $A$. This is again a valid colouring of the rows in
$R$, with $a_j$ replaced by $\tau(a_j) \ne a_j$. If row $i_0$ fails again, then $\tau(a_j) \in
L(i_0, j) \setminus \{a_j\} = B$ for all $j$. So $\tau$ maps two distinct colours of $A \setminus B$
into $A \cap B$, and $|A \cap B| \ge 2$, which forces $B \subseteq A$ and $|A \setminus B| = 1$, a
contradiction.
\end{proof}

\begin{theorem}[Odd strip lemma for $m = 5$; computer-assisted, W7]\label{thm:strip5}
For each of the $74$ end weights $\mu$ of Computation~\ref{comp:windows}(b), in either
orientation of the strip, the coefficient $c_\mu(\ell)$ of $C_5 \T P_\ell$ is nonzero for
\emph{every} odd $\ell \ge 3$. Hence the strip can be coloured with the pin and the forbidden
colours of any of these shapes, at every odd length.
\end{theorem}

\begin{proof}[Proof sketch]
$c_\mu(\ell) = u^{\mathsf T} T^{\ell - 2} v$ with the $30 \times 30$ transfer matrix $T$ of
Theorem~\ref{thm:edge}, whose minimal polynomial is $x\, h_1(x^2)\, h_2(x^2)$ with $h_1 = y^2 +
10y + 5$ and $h_2 = y^4 + 35y^3 + 110y^2 + 100y + 25$, both irreducible over $\mathbb{Q}$ with
only negative real roots. So $s_j = c_\mu(2j+1)$ satisfies the recurrence $h_1 h_2$ for $j \ge 1$,
and its minimal recurrence over $\mathbb{Q}$ is $h_1$, $h_2$ or $h_1h_2$ (found exactly from
$s_1, \dots, s_{13}$; this suffices since $13 - 6 \ge 6$). It is $h_2$ for $96$ and $h_1$ for $52$
of the $148$ cases. Minimality over $\mathbb{Q}$ is minimality over $\mathbb{C}$, so every root has
a nonzero amplitude. The root of largest modulus is unique, and certified ball arithmetic shows
that its term outweighs all others already at $j = 1$, hence for all $j \ge 1$. The predicted
recurrence was also confirmed on exact values up to $\ell = 41$. (Script
\texttt{ws-W7/prove\_strip5.py}.)
\end{proof}

This settles item A1 of the work plan for $m = 5$ (the midpoint rule's nonvanishing for all odd
lengths) and more. Putting the pieces together:

\begin{theorem}[Reduction for $m = 5$]\label{thm:red5}
Let $n \ge 5$ be odd and $L$ a 3-list assignment of $C_5 \T C_n$. Then $C_5 \T C_n$ is
$L$-colourable if $L$ is row-constant, column-constant, has one free row, or has a good cut.
\end{theorem}

\begin{proof}
The first two cases are Corollary~\ref{cor:rowconst}(b), the third is Lemma~\ref{lem:freerow}.
A good cut at a trigger $(r, j)$ uses a shape whose strip coefficient is nonzero at length
$n - 2$ (Theorem~\ref{thm:strip5}); the strip colouring it produces is one of the adversary's
options, and goodness says the two deleted columns can then be filled.
\end{proof}

So the whole of $C_5 \T C_n$ rests on one statement.

\begin{problem}[Window statement]\label{prob:ext}
Every 3-list assignment of $C_5 \T C_n$ ($n$ odd) in which every cut is bad is row-constant,
column-constant, or has one free row.
\end{problem}

By Computation~\ref{comp:windows}(c) this holds for all assignments from $4$ colours and
$n \in \{5, 7\}$, so Theorem~\ref{thm:red5} proves those cases. With more colours the window
picture is richer: changing a single list of a bad $4$-colour window to a 3-subset of $6$ colours
gives $992$ further bad windows (with $\alpha$ fresh), for example
$L_1 = (\{0,1,3\}^3, \{0,3,4\}, \{0,1,3\})$, $L_2 = (\{0,1,3\}^3, \{0,2,3\}, \{0,1,3\})$. Evaluated for each
admissible $\alpha$ separately ($4\,728$ bad pairs of window and $\alpha$), these windows still
produce no new obstruction: for $n = 5$ the assignments in which every cut is bad are again only
column-constant ($24$) or one free row ($39\,120$, counting cyclic shifts). This covers only the
windows within one changed list of a bad $4$-colour window. An exhaustive check of all
$6.7 \times 10^8$ window types with at most $5$ colours is about $370$ CPU-hours (a day on an 18-core machine); a proof for all
palettes needs an argument, not enumeration.

%=====================================================================================
\section{Verification}\label{sec:verification}

All paths are relative to \texttt{odd-torus/}.

\begin{center}\small
\begin{tabular}{@{}p{0.36\textwidth}p{0.3\textwidth}p{0.26\textwidth}@{}}
\toprule
Claim & Script & Output\\
\midrule
Strip coefficient closed form, $\ell$ from 2 to 11 & \texttt{ws-W1/strip\_coeff.py}, \texttt{direct\_expand.py}, \texttt{compare.py} & 0 mismatches in 500\\
Lemma~\ref{lem:prism} & \texttt{ws-W9/prism\_sizes.c} & 8\,335 non-colourable, 0 exceptions\\
Prism lemma (Z)/(W) & \texttt{ws-W9/prism.c} & 986 non-colourable, 0 mismatches\\
Case~1 extension, palettes 5, 6 & \texttt{ws-W9/ext3.c} & 0 failures\\
Control (no $\beta$) & \texttt{ws-W9/ext3\_nododge.c} & 14\,592 failures\\
Theorem~\ref{thm:edge} & \texttt{ws-W4/} & see \texttt{ws-W4/REPORT.md}\\
Computation~\ref{comp:c5c5} & \texttt{ws-W2/certify.py} & \texttt{certify\_t55q4\_nosym.txt}\\
Midpoint rule, $m = 5$ & \texttt{ws-W9/midpoint.py} & 0 mismatches in 2\,600\\
Windows, plan shape, $m = 5$ & \texttt{ws-W8/window5.c} & 6\,904 bad of 1\,048\,576\\
Windows, all 74 shapes & \texttt{ws-W8/window5g.c}, \texttt{shapes5.py} & 400; 184 with real $\alpha$\\
Window code cross-check & \texttt{ws-W8/validate\_window.py} & agrees on 6 windows\\
Lemma~\ref{lem:freerow} construction & \texttt{ws-W9/freerow\_check.py} & 0 failures in 20\,000\\
Theorem~\ref{thm:strip5} & \texttt{ws-W7/prove\_strip5.py} & 148 cases proved\\
Bad windows, 5--6 colours & \texttt{ws-W8/nbhd.py}, \texttt{window5}, \texttt{window5g} & 992 (one list changed)\\
Every cut bad, 5--6 colours, $n = 5$ & \texttt{ws-W8/global\_alpha.py} & only column-constant, one free row\\
Every cut bad, $n = 5, 7$ & \texttt{ws-W8/global\_q4\_all\_n.py} & 600 and 2\,520, all one free row\\
Kernel orientation counts & \texttt{ws-W5/} & see \texttt{ws-W5}\\
\bottomrule
\end{tabular}
\end{center}

%=====================================================================================
\begin{thebibliography}{LSPG}
\bibitem[A]{A} N.~Alon, Combinatorial Nullstellensatz, \emph{Combin. Probab. Comput.} 8 (1999),
7--29.
\bibitem[CWZ]{CWZ} L.~Cai, W.~Wang and X.~Zhu, Choosability of toroidal graphs without short
cycles, \emph{J. Graph Theory} 65 (2010), 1--15.
\bibitem[ERT]{ERT} P.~Erd\H{o}s, A.~L.~Rubin and H.~Taylor, Choosability in graphs,
\emph{Congr. Numer.} 26 (1979), 125--157.
\bibitem[G]{G} F.~Galvin, The list chromatic index of a bipartite multigraph, \emph{J. Combin.
Theory Ser. B} 63 (1995), 153--158.
\bibitem[KP]{KP} R.~N.~Karasev and F.~V.~Petrov, Partitions of nonzero elements of a finite field
into pairs, \emph{Israel J. Math.} 192 (2012), 143--156.
\bibitem[La]{La} M.~Laso\'n, A generalization of Combinatorial Nullstellensatz,
\emph{Electron. J. Combin.} 17 (2010), \#N32.
\bibitem[LSPG]{LSPG} Z.~Li, Z.~Shao, F.~Petrov and A.~Gordeev, The Alon--Tarsi number of a
toroidal grid, \emph{European J. Combin.} 111 (2023), 103697.
\bibitem[S]{S} List-colouring and chromatic-choosability: a dynamic survey, arXiv:2606.31702.
\end{thebibliography}

\end{document}
